% OWNER: Member 1. DRAFT — replace every \todo{} before submission.
\section{Introduction}

Reinforcement learning agents are usually trained and tested on environments
whose reward function is known in advance. ARC-AGI-3~\citep{arcagi3} removes
that assumption: each game is a small grid world whose controls, rules and goal
are not stated, so an agent must explore, infer what counts as progress and
exploit it within a limited number of actions. Performance is measured by
Relative Human Action Efficiency (RHAE), which rewards completing levels with
as few actions as a first-time human player.

We study how far two standard deep RL algorithms can go in this setting when
they learn \emph{online}, from scratch, inside each game: Deep Q-Networks
(DQN; \citealp{mnih2015dqn}), covered in the course, and Proximal Policy
Optimisation (PPO; \citealp{schulman2017ppo}), which is not. Both share one
observation encoder, action space and evaluation protocol, so differences are
attributable to the learning algorithm and to the improvements we add.

\paragraph{Research question.} \todo{One sentence, e.g.: Which mechanisms
--- exploration bonuses, memory, or structured action spaces --- most improve
the sample efficiency of value-based and policy-gradient agents on unseen
ARC-AGI-3 games?}

\paragraph{Contributions.}
\begin{itemize}
  \item A documented Gymnasium-style adapter for the official ARC-AGI-3
        toolkit, with action masking and explicit reset and budget semantics
        (Section~\ref{sec:problem}).
  \item A reproducible evaluation protocol with matched interaction budgets,
        three seeds, and a frozen development / held-out game split
        (Section~\ref{sec:setup}).
  \item Two evidence-driven improvement steps for each algorithm, with
        ablations (Sections~\ref{sec:dqn}--\ref{sec:ppo}).
  \item \todo{Headline result in one line, with official score and Kaggle
        submission ID.}
\end{itemize}
